\documentclass{beamer}
\usepackage[latin1]{inputenc}
\beamertemplateshadingbackground{red!10}{blue!10}
%\usepackage{fancybox}
\usepackage{epsfig}
\usepackage{verbatim}
\usepackage{url}
%\usepackage{graphics}
%\usepackage{xcolor}
\usepackage{fancybox}
\usepackage{moreverb}
%\usepackage[all]{xy}
\usepackage{listings}
\usepackage{filecontents}
\usepackage{graphicx}

\lstset{
  language=Lisp,
  basicstyle=\scriptsize\ttfamily,
  keywordstyle={},
  commentstyle={},
  stringstyle={}}

\def\inputfig#1{\input #1}
\def\inputeps#1{\includegraphics{#1}}
\def\inputtex#1{\input #1}

\inputtex{logos.tex}

%\definecolor{ORANGE}{named}{Orange}

\definecolor{GREEN}{rgb}{0,0.8,0}
\definecolor{YELLOW}{rgb}{1,1,0}
\definecolor{ORANGE}{rgb}{1,0.647,0}
\definecolor{PURPLE}{rgb}{0.627,0.126,0.941}
\definecolor{PURPLE}{named}{purple}
\definecolor{PINK}{rgb}{1,0.412,0.706}
\definecolor{WHEAT}{rgb}{1,0.8,0.6}
\definecolor{BLUE}{rgb}{0,0,1}
\definecolor{GRAY}{named}{gray}
\definecolor{CYAN}{named}{cyan}

\newcommand{\orchid}[1]{\textcolor{Orchid}{#1}}
\newcommand{\defun}[1]{\orchid{#1}}

\newcommand{\BROWN}[1]{\textcolor{BROWN}{#1}}
\newcommand{\RED}[1]{\textcolor{red}{#1}}
\newcommand{\YELLOW}[1]{\textcolor{YELLOW}{#1}}
\newcommand{\PINK}[1]{\textcolor{PINK}{#1}}
\newcommand{\WHEAT}[1]{\textcolor{wheat}{#1}}
\newcommand{\GREEN}[1]{\textcolor{GREEN}{#1}}
\newcommand{\PURPLE}[1]{\textcolor{PURPLE}{#1}}
\newcommand{\BLACK}[1]{\textcolor{black}{#1}}
\newcommand{\WHITE}[1]{\textcolor{WHITE}{#1}}
\newcommand{\MAGENTA}[1]{\textcolor{MAGENTA}{#1}}
\newcommand{\ORANGE}[1]{\textcolor{ORANGE}{#1}}
\newcommand{\BLUE}[1]{\textcolor{BLUE}{#1}}
\newcommand{\GRAY}[1]{\textcolor{gray}{#1}}
\newcommand{\CYAN}[1]{\textcolor{cyan }{#1}}

\newcommand{\reference}[2]{\textcolor{PINK}{[#1~#2]}}
%\newcommand{\vect}[1]{\stackrel{\rightarrow}{#1}}

% Use some nice templates
\beamertemplatetransparentcovereddynamic

\newcommand{\A}{{\mathbb A}}
\newcommand{\degr}{\mathrm{deg}}

\title{Bootstrapping \commonlisp{} on \commonlisp{}}

\author{Ir�ne Durand\\Robert Strandh}
\institute{
LaBRI, University of Bordeaux
}
\date{April, 2019}

%\inputtex{macros.tex}


\begin{document}
\frame{
\resizebox{3cm}{!}{\includegraphics{Logobx.pdf}}
\hfill
\resizebox{1.5cm}{!}{\includegraphics{labri-logo.pdf}}
\titlepage
\vfill
\small{European Lisp Symposium, Genoa, Italy \hfill ELS2019}
}

\setbeamertemplate{footline}{
\vspace{-1em}
\hspace*{1ex}{~} \GRAY{\insertframenumber/\inserttotalframenumber}
}

\frame{
\frametitle{Context: The \sicl{} project}

https://github.com/robert-strandh/SICL

Several objectives:

\begin{itemize}
\item Create high-quality \emph{modules} for implementors of
  \commonlisp{} systems.
\item Improve existing techniques with respect to algorithms and data
  structures where possible.
\item Improve readability and maintainability of code.
\item Improve documentation.
\item Ultimately, create a new implementation based on these modules.
\end{itemize}
}


\frame{
\frametitle{How \commonlisp{} systems are built}

Two main categories of systems:

\begin{enumerate}
\item Systems written mainly in a language other than \commonlisp{},
  typically C or C++.
\item Systems written mainly in \commonlisp{}.
\end{enumerate}

According to Rhodes (SBCL: A Sanely-Bootstrappable Common Lisp):

\begin{enumerate}
\item GCL, ECL, CLISP, ABCL (Java), xcl.
\item Allegro, LispWorks, CMUCL, Scieneer CL, CCL, SBCL.
\end{enumerate}


}

\frame{
\frametitle{How \commonlisp{} systems are built}

Of the systems mainly written in \commonlisp{} there are two ways of
making them evolve:

\begin{enumerate}
\item Image-based techniques, i.e., the system evolves by careful
  evolution of an in-memory image, that is then saved to secondary
  memory to become a new executable.
\item The system is built from a set of source files, using an
  existing \commonlisp{} implementation to create an executable for
  the new system.
\end{enumerate}

Again, according to Rhodes, only SBCL uses the second technique,
though CCL is also built from source, even though at the moment, only
an old CCL version can be used for the bootstrapping process.

}

\frame{
\frametitle{Complications with building Common Lisp systems}

Question: Is the build process for creating a \commonlisp{} system
intrinsically harder than the build process for creating a mere
file-translating compiler, even for a language as complex as
\commonlisp{}, and if so, why?

}

\frame{
\frametitle{Complications with building Common Lisp systems}

Possible answer: We think that it is, and we think that it is due to
the fact that:

\begin{itemize}
\item a file translating compiler contains essentially only code, and
  code is generated from source statements and expressions as the
  implementation language makes available, whereas
\item a \commonlisp{} system also contains complex data structures,
  and these data structures are best created programmatically.
\end{itemize}

Imagine, for instance, writing the C source statements that will
generate a single generic function with its methods, effective method
functions, etc.

}

\frame{
\frametitle{Strategies for creating the data structures}

There are several possible strategies for creating these data
structures:

\begin{enumerate}
\item Create them when the system starts.  This strategy is used by
  ECL and Clasp.
\item Create them in an executing subset \commonlisp{} system, and
  then save the memory image to a file.  This strategy is used by the
  image-based implementations, but also by SBCL.
\item Create a mirror of the final executable as an octet vector in
  the memory of a host \commonlisp{} system.  This strategy is used by
  SBCL to create an initial ``core'' image.  It is also the strategy
  used by SICL.
\end{enumerate}


}

\frame{
\frametitle{Previous work}


}

\frame[containsverbatim]{
\frametitle{Definition of the class \texttt{t}}

\begin{verbatim}
(defclass t ()
  ()
  (:metaclass built-in-class))
\end{verbatim}

}

\frame{
\frametitle{Future work}


}

\frame{
  \frametitle{Acknowledgments}

}

\frame{
\frametitle{Thank you}

Questions?
}

%% \frame{\tableofcontents}
%% \bibliography{references}
%% \bibliographystyle{alpha}

\end{document}
